\chapter{Money, the Price Level, and Inflation (Text Ch. 24)}

\section{What is money?}

\begin{definition}[Money]
\emph{Money} is any commodity or token that is generally accepted as a \emph{means of
payment}. A means of payment is a method of settling a debt.
\end{definition}

Money has three functions.
\begin{enumerate}
  \item \textbf{Medium of exchange}: an object generally accepted in exchange for goods
        and services. Without money, people would have to \emph{barter} (trade goods
        directly for goods), which requires a \emph{double coincidence of wants}: you must
        find someone who has what you want and wants what you have. Money removes this
        requirement.
  \item \textbf{Unit of account}: an agreed measure for stating the prices of goods and
        services. With $n$ goods and no unit of account we would need
        $n(n-1)/2$ relative prices; with money we need only $n$ prices.
  \item \textbf{Store of value}: something that can be held and exchanged later for goods
        and services. Money is a good (not perfect) store of value: its value falls with
        inflation.
\end{enumerate}

\subsection{Money in Canada today}

Money in Canada consists of
\begin{itemize}
  \item \textbf{currency}: the notes (issued by the Bank of Canada) and coins (from the
        Royal Canadian Mint) held by individuals and businesses (not those held by banks);
        and
  \item \textbf{deposits} at banks and other depository institutions.
\end{itemize}

\subsection{Official measures of money}

\begin{definition}[M1 and M2]
\textbf{M1} consists of currency held by individuals and businesses plus chequable
deposits owned by individuals and businesses. \textbf{M2} consists of M1 plus all other
deposits: non-chequable deposits and fixed term deposits.
\end{definition}

All the items in M1 are means of payment, so they are money. Some savings deposits in M2
are not means of payment: they are \emph{liquid assets} that can be converted into a means
of payment quickly and at a known price. (The Bank of Canada also publishes broader
measures, M1+ and M2+, that include credit unions and trust companies.)

\begin{important}[What is not money]
\begin{itemize}
  \item \textbf{Cheques} are instructions to a bank to transfer money; the
        \emph{deposit} is money, the cheque is not.
  \item \textbf{Credit cards} are not money: they are an ID that lets you take a loan,
        which you later repay with money.
  \item \textbf{Debit cards} and e-transfers are ways of moving deposits; the deposit is
        the money.
\end{itemize}
\end{important}

\section{Depository institutions}

\begin{definition}[Depository institution]
A \emph{depository institution} is a firm that takes deposits from households and firms
and makes loans to other households and firms.
\end{definition}

Types of depository institutions in Canada:
\begin{enumerate}
  \item \textbf{Chartered banks}: private firms chartered under the \emph{Bank Act} of
        1991 to receive deposits and make loans (RBC, TD, Scotiabank, BMO, CIBC, \dots).
  \item \textbf{Credit unions and caisses populaires}: cooperative organizations owned by
        their members, operating under the \emph{Cooperative Credit Associations Act} of 1991
        and provincial laws.
  \item \textbf{Trust and mortgage loan companies}: privately owned institutions that
        receive deposits, make loans, and act as trustees for estates and pension funds.
\end{enumerate}

A bank's balance sheet:
\begin{itemize}
  \item liabilities: deposits (and the owners' equity);
  \item assets: \emph{reserves} (notes and coins in its vault plus its deposit at the
        Bank of Canada), liquid assets (Government of Canada Treasury bills), securities,
        and loans (the highest-yielding but least liquid assets).
\end{itemize}
A bank makes a profit from the spread between the interest it earns on loans and
securities and the interest it pays on deposits. It balances profit against the risk of
not being able to meet depositors' withdrawals.

Depository institutions provide four economic benefits:
\begin{enumerate}
  \item \textbf{create liquidity}: they borrow short (deposits) and lend long (mortgages);
  \item \textbf{pool risk}: a loan default spreads over many depositors;
  \item \textbf{lower the cost of borrowing}: one large loan from a bank replaces finding
        many small lenders;
  \item \textbf{lower the cost of monitoring borrowers}: specialists check
        creditworthiness once for many depositors.
\end{enumerate}

\section{The Bank of Canada}

\begin{definition}[Central bank]
The \emph{Bank of Canada} (established by the Bank of Canada Act, 1934; opened 1935) is
Canada's central bank: a public authority that supervises financial institutions and
markets and conducts monetary policy.
\end{definition}

The Bank of Canada is
\begin{itemize}
  \item the \textbf{banker to banks}: banks hold deposits at the Bank and use them to
        settle payments with each other;
  \item the \textbf{banker to the government}: it holds the government's deposits and
        manages the government's borrowing (selling bonds and Treasury bills);
  \item the \textbf{lender of last resort}: it stands ready to lend reserves to banks that
        are short of them, preventing a liquidity problem from becoming a bank run;
  \item the \textbf{sole issuer of bank notes}.
\end{itemize}

Its balance sheet shows assets (mainly Government of Canada securities, plus loans to
banks) and liabilities (bank notes in circulation, banks' deposits (``settlement
balances''), and government deposits).

\begin{definition}[Monetary base]
The \emph{monetary base} is the sum of Bank of Canada notes, coins and banks' deposits at
the Bank of Canada.
\end{definition}

\textbf{Policy tools.} The Bank of Canada influences the quantity of money and interest
rates using
\begin{enumerate}
  \item the \textbf{overnight rate target} (its main instrument; see the monetary policy
        chapter): the interest rate on overnight loans between banks;
  \item \textbf{open market operations}: buying or selling government securities in the
        open market. Buying securities pays with new reserves and \emph{increases} the
        monetary base; selling decreases it;
  \item the \textbf{bank rate}: the rate at which it lends reserves to banks (set at the
        top of the operating band, 0.25 points above the target);
  \item the \textbf{settlement balances rate}: the rate it pays on banks' reserves (at
        the bottom of the band).
\end{enumerate}
Canada has had no legally required reserve ratio since 1994; banks hold the reserves
they \emph{desire}.

\section{How banks create money}

Banks create money when they make loans: a loan is deposited in someone's account, and
deposits are money. Banks create money by lending their \emph{excess reserves}
(actual reserves minus desired reserves).

\begin{example}[Deposit creation with no currency drain]
Banks want reserves equal to $b = 10\%$ of deposits and the public holds no extra
currency. A new deposit of \$1{,}000 gives the bank \$900 of excess reserves, which it
lends. The \$900 is spent and redeposited, and that bank lends \$810, and so on:
\[
  \Delta D = 1000 + 900 + 810 + \cdots = 1000 \cdot \frac{1}{1 - 0.9} = \frac{1000}{0.1}
  = \$10{,}000.
\]
\end{example}

The quantity of deposits that banks can create is limited by three factors.
\begin{enumerate}
  \item \textbf{The monetary base}: its size limits the total quantity of money the
        banking system can create, because banks have desired reserves and households and
        firms have desired currency holdings, and both must come out of the base.
  \item \textbf{Desired reserves}: the reserves (notes and coins in its vault plus its
        deposit at the Bank of Canada) that a bank plans to hold. The \emph{desired
        reserve ratio} is the ratio of desired reserves to deposits.
  \item \textbf{Desired currency holding}: people hold some money as currency, so part of
        each loan leaks out of the banking system (a \emph{currency drain}). The
        \emph{currency drain ratio} is the ratio of currency to deposits.
\end{enumerate}

\subsection{The money multiplier}

\begin{definition}[Money multiplier]
The \emph{money multiplier} is the ratio of the change in the quantity of money to the
change in the monetary base.
\end{definition}

\begin{theorem}[Money multiplier formula]
Let $a = C/D$ be the currency drain ratio and $b = R/D$ the desired reserve ratio, where
$C$ is currency, $D$ deposits and $R$ reserves. Then
\[
  \text{money multiplier} = \frac{M}{MB} = \frac{1 + a}{a + b}.
\]
\end{theorem}

\begin{proof}
Money is $M = C + D = aD + D = (1 + a)D$. The monetary base is held either as currency
or as reserves: $MB = C + R = aD + bD = (a + b)D$. Dividing,
$M/MB = (1 + a)/(a + b)$. The same ratio holds for changes.
\end{proof}

\begin{example}
With $a = 0.5$ and $b = 0.1$ the multiplier is $1.5/0.6 = 2.5$. If the Bank of Canada
buys \$100 million of securities, the monetary base rises by \$100 million and the
quantity of money rises by \$250 million. Deposits rise by $100/0.6 \approx 166.7$,
currency by $0.5 \times 166.7 \approx 83.3$ and reserves by $0.1 \times 166.7 \approx 16.7$;
check: $83.3 + 16.7 = 100$ (the base) and $83.3 + 166.7 = 250$ (money).
\end{example}

\begin{important}[Pitfall]
With no currency drain ($a = 0$) the formula reduces to the simple deposit multiplier
$1/b$. A larger currency drain or a larger desired reserve ratio makes the multiplier
\emph{smaller}.
\end{important}

\section{The money market}

\subsection{The influences on money holding}

The quantity of money that people plan to hold depends on four main factors.
\begin{enumerate}
  \item \textbf{The price level}: the quantity of \emph{nominal} money demanded is
        proportional to the price level (if prices double, people hold twice as many
        dollars). The quantity of \emph{real} money, $M/P$, is independent of the price
        level.
  \item \textbf{The nominal interest rate}: it is the opportunity cost of holding money
        (instead of interest-bearing bonds). A higher interest rate decreases the
        quantity of real money demanded.
  \item \textbf{Real GDP}: higher real GDP means more spending, so more money is held.
  \item \textbf{Financial innovation}: e.g.\ daily-interest chequing, debit cards, online
        banking and e-transfers reduce the demand for currency and for money in general.
\end{enumerate}

\begin{definition}[Demand for money]
The \emph{demand for money} (MD) is the relationship between the quantity of real money
demanded and the nominal interest rate, other things remaining the same. It slopes
downward. A rise in real GDP shifts it rightward; financial innovation generally shifts it
leftward.
\end{definition}

\subsection{Money market equilibrium}

In the \emph{short run} the price level is fixed and the Bank of Canada determines the
quantity of money, so the supply of real money (MS) is vertical. The nominal interest
rate adjusts to make the quantity demanded equal the quantity supplied. If the interest
rate is above equilibrium, people hold more money than they want and buy bonds; bond prices
rise and the interest rate falls.

\begin{center}
\begin{tikzpicture}[scale=0.75]
  \draw[->] (0,0) -- (8,0) node[right] {Real money (\$ billions of 2002)};
  \draw[->] (0,0) -- (0,8) node[above] {Nominal interest rate (\%)};
  \draw[thick, blue] (1,7) -- (7,1) node[right] {MD};
  \draw[thick, red] (4,0.5) -- (4,7.5) node[above] {MS$_0$};
  \draw[thick, red, dashed] (5,0.5) -- (5,7.5) node[above] {MS$_1$};
  \fill (4,4) circle (2.5pt);
  \fill (5,3) circle (2.5pt);
  \draw[dotted] (0,4) node[left] {$i_0$} -- (4,4);
  \draw[dotted] (0,3) node[left] {$i_1$} -- (5,3);
\end{tikzpicture}
\end{center}

An increase in the quantity of money (MS$_0$ to MS$_1$) lowers the short-run equilibrium
interest rate. In the \emph{long run}, real GDP equals potential GDP and the real
interest rate is determined in the loanable funds market; the \emph{price level} adjusts
to make the quantity of real money supplied equal the quantity demanded. An increase in
the quantity of money therefore raises the price level in the same proportion in the long
run.

\section{The quantity theory of money}

\begin{definition}[Velocity and the equation of exchange]
The \emph{velocity of circulation} $V$ is the average number of times a dollar of money
is used annually to buy the goods and services that make up GDP:
\[
  V = \frac{PY}{M},
\]
where $P$ is the price level, $Y$ is real GDP and $PY$ is nominal GDP. Rearranged, this is
the \emph{equation of exchange}:
\[
  MV = PY.
\]
\end{definition}

The equation of exchange is an identity. It becomes a theory under two assumptions:
velocity is not influenced by the quantity of money, and potential GDP is not influenced
by the quantity of money.

\begin{theorem}[Quantity theory of money]
In the long run, an increase in the quantity of money brings an equal percentage increase
in the price level. In growth rates,
\[
  \pi = \mu + v - g,
\]
where $\pi$ is the inflation rate, $\mu$ the money growth rate, $v$ the growth rate of
velocity and $g$ the growth rate of real GDP. With constant velocity,
$\pi = \mu - g$.
\end{theorem}

\begin{proof}[Sketch]
Taking percentage changes of $MV = PY$ (the percentage change of a product is
approximately the sum of the percentage changes): $\mu + v = \pi + g$.
\end{proof}

\begin{example}
$M = \$500$ billion, $V = 4$ and real GDP $Y = \$1{,}000$ billion. Then
$P = MV/Y = 2000/1000 = 2$ (a price index of 200). If money grows at 7\% a year, velocity is
constant and real GDP grows at 3\%, inflation is $\pi = 7 + 0 - 3 = 4\%$ a year.
If money growth is cut to 3\%, long-run inflation falls to 0.
\end{example}

Historical and international evidence supports the quantity theory in the long run:
countries with high money growth (e.g.\ in hyperinflations) have high inflation. In the
short run the correlation is much weaker because velocity and real GDP fluctuate.

\begin{remark}[Inflation and the nominal interest rate]
In the long run, the nominal interest rate equals the real interest rate plus the
inflation rate (the \emph{Fisher effect}): $i \approx r + \pi$. Higher money growth means
higher inflation and higher nominal interest rates.
\end{remark}

\section{Summary}

\begin{itemize}
  \item Money is a generally accepted means of payment; functions: medium of exchange,
        unit of account, store of value. M1 $=$ currency $+$ chequable deposits; M2 $=$ M1
        $+$ other deposits. Credit cards are not money.
  \item Depository institutions (chartered banks, credit unions, trust companies) create
        liquidity, pool risk and lower borrowing and monitoring costs.
  \item The Bank of Canada is banker to banks and government, lender of last resort and
        issuer of notes; monetary base $=$ notes $+$ coins $+$ banks' deposits at the Bank.
  \item Banks create money by lending excess reserves; money multiplier
        $= (1 + a)/(a + b)$.
  \item Demand for real money falls with the nominal interest rate and rises with real
        GDP; short-run money market equilibrium determines the interest rate, long-run
        equilibrium the price level.
  \item Quantity theory: $MV = PY$, $\pi = \mu + v - g$; in the long run money growth
        causes inflation.
\end{itemize}
